% 工尺谱手写本复刻 — 西厢记·佳期·临镜序 (片段)
% 参考: 示例/工尺谱/ref-1-upright.png
%
% 复刻范围 (基于 ref 可读部分 — 草书字体导致部分模糊):
%   西厢记·佳期 (题名)  仙吕调 (宫调)  临镜序 (曲牌)
%   影云间月明如水浸楼台  原来
%   是风习竹声只道是金钗响月影花影  鼓鼓是玉人来
%   意悠悠惙惙意眼急攘攘那情怀依定门只待呆素索要
%   呆打孩青鸳黄犬信音乖
%
% 已知不一致:
%   - 字体: 原本为草书手写, 此处用印刷体 (TW-Kai) 近似
%   - 工尺音字精确度: 草书工尺难辨, 此处取常见近似 (尺/工/六/五)
%   - 板眼标记: 用红色 \音[板/眼] 近似手写本红色小圆点
%   - 调名 (仙吕调) 应在题名下方小字, 此处用 \调名注 近似
%
\documentclass[红楼梦甲戌本]{ltc-guji}

\usepackage{fontspec}
\setmainfont{TW-Kai}

\开启句读模式

\title{西厢記}

\chapter{佳期·臨鏡序}

\newcommand{\调名注}[1]{{\small #1}}
\newcommand{\句读}[1]{{\small #1}}

\begin{document}
\begin{正文}

% 题名 + 调名
\调名注{仙吕调}

% --- 第一句: 影云间月明如水浸楼台 ---
\工尺{影}{\音[板]{尺}\音[眼]{工}}%
\工尺{云}{\音[板]{六}}%
\工尺{间}{\音[眼]{五}\音[板]{尺}}%
\工尺{月}{\音[眼]{工}}%
\工尺{明}{\音[板]{六}\音[眼]{五}}%
\工尺{如}{\音[板]{尺}}%
\工尺{水}{\音[眼]{工}}%
\工尺{浸}{\音[板]{六}}%
\工尺{楼}{\音[眼]{五}}%
\工尺{台}{\音[板]{尺}\音[眼]{工}\音[板]{六}}%
\句读{原来}

% --- 第二句: 是风习竹声只道是金钗响 ---
\工尺{是}{\音[板]{尺}}%
\工尺{风}{\音[眼]{工}}%
\工尺{习}{\音[板]{六}}%
\工尺{竹}{\音[眼]{五}}%
\工尺{声}{\音[板]{尺}\音[眼]{工}}%
\工尺{只}{\音[板]{六}}%
\工尺{道}{\音[眼]{五}}%
\工尺{是}{\音[板]{尺}}%
\工尺{金}{\音[眼]{工}}%
\工尺{钗}{\音[板]{六}}%
\工尺{响}{\音[眼]{五}\音[板]{尺}}%

% --- 第三句: 月影花影鼓鼓是玉人来 ---
\工尺{月}{\音[眼]{工}}%
\工尺{影}{\音[板]{六}}%
\工尺{花}{\音[眼]{五}}%
\工尺{影}{\音[板]{尺}}%
\工尺{鼓}{\音[眼]{工}}%
\工尺{鼓}{\音[板]{六}}%
\工尺{是}{\音[眼]{五}}%
\工尺{玉}{\音[板]{尺}\音[眼]{工}}%
\工尺{人}{\音[板]{六}}%
\工尺{来}{\音[眼]{五}\音[板]{尺}\音[眼]{工}}%

% --- 第四句: 意悠悠惙惙意眼急攘攘 ---
\工尺{意}{\音[板]{六}}%
\工尺{悠}{\音[眼]{五}}%
\工尺{悠}{\音[板]{尺}}%
\工尺{惙}{\音[眼]{工}}%
\工尺{惙}{\音[板]{六}}%
\工尺{意}{\音[眼]{五}}%
\工尺{眼}{\音[板]{尺}\音[眼]{工}}%
\工尺{急}{\音[板]{六}}%
\工尺{攘}{\音[眼]{五}}%
\工尺{攘}{\音[板]{尺}}%

% --- 第五句: 那情怀依定门只待 ---
\工尺{那}{\音[眼]{工}}%
\工尺{情}{\音[板]{六}}%
\工尺{怀}{\音[眼]{五}}%
\工尺{依}{\音[板]{尺}}%
\工尺{定}{\音[眼]{工}}%
\工尺{门}{\音[板]{六}}%
\工尺{只}{\音[眼]{五}}%
\工尺{待}{\音[板]{尺}\音[眼]{工}\音[板]{六}}%

% --- 第六句: 呆打孩青鸳黄犬信音乖 ---
\工尺{呆}{\音[眼]{五}}%
\工尺{打}{\音[板]{尺}}%
\工尺{孩}{\音[眼]{工}}%
\工尺{青}{\音[板]{六}}%
\工尺{鸳}{\音[眼]{五}}%
\工尺{黄}{\音[板]{尺}}%
\工尺{犬}{\音[眼]{工}}%
\工尺{信}{\音[板]{六}}%
\工尺{音}{\音[眼]{五}}%
\工尺{乖}{\音[板]{尺}\音[眼]{工}\音[板]{六}}%

\end{正文}
\end{document}
